\section*{The Propers: Themes and Movement}
\label{triptych:brief-synthesis:start}
\begingroup\footnotesize
\setlength{\parskip}{0.32em}

\begin{studybox}{Governing thesis}
God's covenant fidelity becomes gift in Christ, and grace creates its own
living answer. The poor plead a promise they did not make; faith receives
Christ's word, hope walks before sight, and charity returns to the Giver in
thanksgiving. At the altar that return enters Christ's propitiatory oblation;
in Communion the promised Seed gives himself as heavenly bread. The Mass ends
by sending the nourished pilgrim forward toward redemption's eternal fullness.
\end{studybox}

\subsection*{1. The covenant's poor call upon the promised Seed\enspace\properrefs{Int., Ep., Grad.}}

The Introit begins from devastation rather than possession. Psalm~73's
sanctuary has been violated, the poor appear abandoned, and God's honor is
mocked in the humiliation of his servants. Yet the people still know how to
pray: \latin{Réspice, Dómine, in testaméntum tuum}. Their first argument is not
their strength or merit, but God's own covenant. Augustine hears here
Jeremiah's new covenant and the heavenly inheritance: ``we possess the tablets
and await the inheritance'' (\emph{Enarr. in Ps.~73} 23). Cassiodorus says that
those with no sufficiency of their own plead the unfailing promise of eternal
life; St~Anthony of Padua sees the testament ratified in Christ's Blood and the
poor soul's inheritance in Christ himself. Literal Israel's lament and the
Church's christological prayer therefore remain joined: the God invoked is the
one who gave the covenant, preserved the people, and brought the promise to
fulfilment.

Galatians names that fulfilment. The promise was spoken to Abraham and to his
one Seed, Christ; the Law, given later, cannot annul God's disposition of the
inheritance. Jerome and Theodoret insist on Christ as the personal center of
the singular; Augustine adds the ecclesial completion supplied by Paul himself:
those incorporated into Christ are gathered into the one seed and made heirs
according to promise (Gal.~3:27--29). Thus the Psalm's poor are not spectators
outside the promise. In Christ they plead from within the gift, holding its
pledge while awaiting its unveiled inheritance.

The Gradual repeats \latin{Réspice} and adds a line from Psalm~88: ``remember
the reproach of thy servants.'' That psalm laments the seemingly defeated
Davidic covenant and immediately speaks of reproach against the footsteps of
God's Anointed. The chant thereby places the servants' suffering inside God's
own cause. Augustine hears Christians bearing reproach because they bear
Christ's name; Cassiodorus sees the despised Christian name itself. The liturgy
does not promise that covenant fidelity removes every historical desolation.
It teaches the poor to carry desolation into the prayer of Christ and to await
the vindication of his kingdom.

\subsection*{2. The good Law serves the promise that gives life\enspace\properrefs{Ep., All., Gosp.}}

Paul's question supplies the boundary: ``Is the Law then against the promises
of God? God forbid.'' Chrysostom calls the Law a bridle restraining sin and a
probe that makes the wounded desire a physician. Theodoret emphasizes its
providential guardianship of the people from whom the Seed would arise.
Augustine sees the command exposing pride so that salvation is sought from the
Mediator rather than from one's own hand. Aquinas gathers these senses: the
Law represses wrongdoing, disciplines desire, manifests weakness, and figures
the grace to come. What it cannot do by external prescription alone is
communicate the Spirit who heals and vivifies. The Law is holy and
pedagogical; promise gives the inheritance.

The Alleluia places Moses within this unity. The traditional title makes
Psalm~89 his prayer, and Moses sings that the Lord has been \emph{our} refuge
through every generation. Augustine hears the minister of the first covenant
beholding the second and Christ joining both. Cassiodorus's Roman lemma,
\latin{a generatióne et progénie}, makes living continuity explicit: mortal
generations proceed from one another while the eternal refuge remains. The
Epistle's \latin{Deus autem unus est} excludes any division between the God of
Abraham, Moses, and Christ. Moses truly ministers; angels truly serve; Christ
alone is the Seed in whom the blessing is given and the Mediator who brings
his members to the inheritance.

The Gospel enacts this order. Jesus sends the ten to the priests under the
Law, but cleanses by his word while they go. Tertullian answers Marcion:
Christ honors the Creator's Law while revealing himself as the true Priest
and temple. Penance retains the order: Christ principally forgives; the
ordained priest, as his minister, judicially absolves and reconciles. Grace
gives command and ministry efficacy dependent on Christ.

\clearpage

\subsection*{3. Faith receives, hope walks, charity returns\enspace\properrefs{Coll., Gosp.}}

The Collect gives the interior form of the journey. Its paired imperatives
ask God to \emph{give} an increase of faith, hope, and charity and to
\emph{make} the petitioners love what he commands. Promise and precept are not
rivals: faith cleaves to the promised Christ, hope tends toward the inheritance,
and charity unites the will to the end and therefore loves the way God commands.
Augustine's \latin{Da quod iubes, et iube quod vis} illuminates this grammar.
The love is truly the creature's voluntary act, yet even that act is asked from
God as grace. Trent accordingly quotes the virtue-increase petition when
teaching growth in received justification, and locates all Christian merit
within Christ's preceding and accompanying gift.

Luke turns those virtues into verbs. Ten sufferers stand at a distance and
lift one voice: \latin{Iesu praecéptor, miserére nostri}. They receive a command
whose effect is not yet visible and set out; only \latin{dum irent} are they
cleansed. Bonaventure reads the final saving word through the three virtues:
faith begins the movement, hope advances it, and charity perfects it. The
Samaritan sees the gift, turns back, glorifies God with a great voice, falls at
Jesus' feet, and gives thanks. Bede therefore distinguishes benefit from its
living reception: all ten are cleansed, but Luke follows one faith into worship,
thanksgiving, and the further word, ``thy faith hath saved thee.''

The Fathers and medieval doctors discover the Church's penitential road in
\latin{ite---osténdite---dum irent}. Augustine, St~Anthony, Albert, Bonaventure,
and their successors contemplate contrition, confession, satisfaction, and
the keys: the heart turns under grace, intends obedient confession, and is
restored to visible communion. Their governing axiom remains Albert's:
\latin{Deus enim solus emundat peccata}, God alone cleanses sins. Yet the return
is no private interior event. Christ heals through an ecclesial economy and
receives the reconciled sinner at his feet. What began as the raised voice of
need becomes the great voice of praise.

\subsection*{4. Mercy received returns as oblation and Eucharist\enspace\properrefs{Gosp., Off., Sec., Comm.}}

The Offertory gives the Samaritan's return a form for the whole assembly:
\latin{In te sperávi, Dómine; dixi: Tu es Deus meus; in mánibus tuis témpora
mea}. Theodoret reads David surrendering the changing distributions of peace
and conflict, grief and joy, to providence; Augustine hears Christ the Head and
his members speaking as one; Cassiodorus sees Christ's assumed humanity and
the hour held in the Father's power. As bread and wine are brought forward,
the faithful place labor, suffering, gratitude, and their final hour within
Christ's self-offering. Hope is not waiting without action: it is the offering
of one's times to the God who holds them.

The Secret then prays twice \latin{propitiáre}: be propitious to thy people;
be propitious to their gifts. Giver and offering both remain beneath mercy.
Tertullian calls the healed Samaritan's thanksgiving the true oblation made to
Christ, the true High Priest and temple. The cure is not bought by gratitude;
gratitude is grace recognizing its source. Positively, the Mass is a true
propitiatory sacrifice because the one sacrifice of Calvary is made
sacramentally present: the same Christ offers through the ministry of priests,
and the fruits of his self-offering are applied for pardon and grace (Trent,
Session~XXII, ch.~2; CCC~1367). Aquinas distinguishes its sacrificial effect as
offered from its sacramental effect as received without dividing the one
Eucharistic mystery.

At Communion the remembered manna becomes present thanksgiving: Wisdom says
that God gave bread to Israel; the Church sings, \latin{Panem de caelo dedísti
nobis, Dómine}. Origen and the two Gregories see the one divine Word nourishing
souls according to their capacity; Rabanus leads the manna through John~6 to
Christ's life-giving flesh and blood. Albert calls the Eucharistic bread
heavenly, life-giving, unitive, viaticum, and pledge of fruition. Aquinas says
manna chiefly figures the sacrament's effect: the one Christ refreshes the mind
in every spiritual good. Bonaventure preserves both letter and mystery---the
historical food truly serves Israel, and in sacramental fulfilment Christ bears
fruit according to grace-enlarged devotion. The stranger who returned is thus
gathered into the grateful plural \latin{nobis}: the gift becomes communion.

\subsection*{5. Heavenly food increases the pilgrim toward eternal refuge\enspace\properrefs{Coll., All., Gosp., Comm., Postcomm.}}

The repeated \latin{augméntum} binds the first prayer to the last. At the
beginning God increases the theological virtues; after Communion the recipients
ask that \emph{we} may advance toward the increase of eternal redemption. The
objective redemption accomplished in Christ needs no enlargement. Its grace,
however, is received ever more deeply by pilgrims. Aquinas compares the
Eucharist to food that sustains, increases, repairs, and delights spiritual
life; actual reception augments grace and charity, while viaticum gives the
power to reach the glory not yet seen. The \emph{Roman Catechism} and
St~John Paul~II make the same passage from Wisdom's every-delight manna to the
Eucharistic food of the journey and foretaste of messianic fullness.

The Gospel ends with a command rather than repose: \latin{Surge, vade}. Bede,
Bonaventure, Albert, and Denis understand the rising as passage from sin into
strong works and the going as advance toward neighbor, God, and the heavenly
homeland. St~Anthony reads \latin{surge} as unity rising to upright works and
\latin{vade} as humility's keeper going securely; his surrounding exposition
also treats return, prostration, thanksgiving, and saving faith. The Mass has traced the same pilgrimage. Faith came
crying, hope walked before sight, charity returned in thanks; sacrifice placed
people and gifts beneath mercy; sacrament gave the promised Christ as food.
Now the Church rises and goes. The refuge of every passing generation draws
her beyond changing times, ruined sanctuaries, and incomplete healing toward
the inheritance where covenant is unveiled, thanksgiving is perpetual, and
redemption's fruit can no longer be lost.

\endgroup

\label{triptych:brief-synthesis:end}
\clearpage
